% !TEX program = xelatex
\documentclass[11pt,a4paper,sans]{moderncv}
\moderncvstyle{classic}
\moderncvcolor{blue}
\moderncvicons{marvosym}

\usepackage[UTF8]{ctex}
\usepackage[top=1.25cm,bottom=1.25cm,left=1.7cm,right=1.7cm]{geometry}
\usepackage{enumitem}
\usepackage{setspace}

\setlength{\hintscolumnwidth}{2.35cm}
\setstretch{1.12}
\setlist[itemize]{label=\textbullet,leftmargin=*,topsep=2pt,partopsep=0pt,parsep=0pt,itemsep=3.5pt}
\nopagenumbers{}

\definecolor{linkcolor}{RGB}{0,92,153}
\AtBeginDocument{\hypersetup{
  urlcolor=linkcolor,
  linkcolor=linkcolor,
  pdfauthor={Max Young (Le Yang)},
  pdftitle={Senior Software Engineer - Autonomous Driving SOTA/FOTA and Fleet Platforms}
}}

% Identity
\firstname{Max}
\familyname{Young \normalfont{{\fontsize{11}{13}\selectfont (Le Yang)}}}
\title{Senior Software Engineer | Autonomous Driving SOTA/FOTA \& Fleet Platforms}
\address{Beijing, China}{}{}
\email{alvysurvey@gmail.com}
\mobile{+86 186 0103 6905}
\photo[58pt][0pt]{../../yangle.jpg}
\extrainfo{%
  \href{https://maxyoung.fun}{Portfolio: maxyoung.fun}\quad
  \href{https://www.linkedin.com/in/max-young-16b690143/}{LinkedIn: Max Young}%
}

\begin{document}
\makecvtitle
\vspace{-0.9cm}

\section{Profile}
\cvitem{}{Senior software engineer with 13+ years of experience building production systems, including 4+ years in autonomous driving. Leads vehicle-platform work from product definition and architecture through implementation and fleet operations. Reduced average SOTA deployment time by 67\% while expanding capacity 25$\times$ from 200 to 5,000 vehicles, and led an autonomous-vehicle FOTA platform to M3/MVP. Strong in Python, React and Linux, with practical C++ experience; open to relocation across Germany.}

\section{Professional Experience}
\cventry{2022--Present}{Software Engineer}{\href{https://rino.ai/}{Rino.ai}}{Beijing, China}{Autonomous Driving}{%
  \begin{itemize}
    \item Led architecture, product definition and technical direction for a three-person vehicle-platform team, turning separate engineering tools into reusable delivery infrastructure for R\&D, factory and fleet operations.
    \item \textbf{SOTA Platform:} Reduced average software deployment time from approximately 30 to 10 minutes (67\%) while increasing maximum capacity from 200 to 5,000 vehicles (25$\times$).
    \item \textbf{FOTA Platform:} Led design and implementation through M3/MVP, covering distributed execution across TBOX/ADCU nodes, upgrade orchestration, failure recovery, power-loss safety, vehicle clients, backend deployment and web-based management.
    \item Unified delivery workflows for software, maps, vehicle configuration and calibration, including vehicle/series targeting, execution verification, service-health monitoring, interrupted-transfer recovery and version/state consistency.
    \item \textbf{Large-map Distribution:} Led architecture trade-offs for hundreds-of-gigabytes datasets under limited disk space and unreliable networks, evaluating rsync, incremental delivery and OverlayFS for storage efficiency, atomic switching, rollback, progress and ETA accuracy.
    \item \textbf{Reliability \& Observability:} Improved fleet-update transparency by analysing rsync transfer-state semantics, interrupted-transfer behaviour and ETA accuracy, and by exploring lightweight health monitoring for vehicle-hosted services across several thousand vehicles.
    \item \textbf{Vehicle Calibration Platform:} Created an end-to-end Python, Shell, C++ and OpenCV workflow used on hundreds of vehicles; increased factory throughput from 2--3 to 10+ vehicles per day and cut field calibration from 0.5--1 day to under 30 minutes.
    \item \textbf{Fleet Operations Platform:} Built a Flask/React/Ant Design system that standardised incident management and route delivery across R\&D and operations; helped reduce the operational incident rate from 13 to 7 per 10,000 km within six months and supported delivery of dozens of new routes.
    \item \textbf{Engineering Infrastructure:} Evaluated GitLab's PyPI Registry for reusable internal Python packages and improved deployment automation using systemd timers, controlled long-running rsync processes, distributed FRP and cloud-hosted backend/web services.
    \item \textbf{Productisation:} Defined the product positioning, web information architecture and maintainable documentation approach for a unified vehicle-management platform, turning project-specific tooling into reusable team and product capabilities.
  \end{itemize}}

\newpage

\section{Professional Experience (continued)}

\cventry{2017--2022}{Software Engineer}{\href{https://dedao.cn/}{Dedao}}{Beijing, China}{Online Education \& E-commerce}{%
  \begin{itemize}
    \item Designed backend systems for inventory, orders and distribution; automated order ingestion and settlement across JD.com, Taobao and Youzan using Python/Django and Go/Gin.
    \item Built financial recognition, settlement and analytics services in Flask, replacing manual office workflows and improving finance-team efficiency.
    \item Introduced Celery/RabbitMQ and Faust/Kafka for asynchronous processing, built a GraphQL gateway for service integration, and championed TDD and contract testing to reduce defects and integration cost.
  \end{itemize}}

\par\vspace{0.32cm}

\cventry{2014--2017}{Software Engineer}{\href{https://kaoshixing.com/}{Kaoshixing}}{Beijing, China}{B2B SaaS}{%
  Led architecture, full-stack development, DevOps and engineering coordination for an online examination platform, taking it from initial build to production scale serving 10,000+ enterprise customers. \textit{Python, Flask, JavaScript, Linux}}

\vspace{0.12cm}

\cventry{2013--2014}{Map Quality Engineer}{\href{https://www.amap.com/}{AutoNavi (Alibaba Group)}}{Beijing, China}{Mapping Services}{%
  Analysed geospatial POI and road-network data, established map-asset quality-control standards, and advised data-collection planning and quality measurement.}

\vspace{0.12cm}

\cventry{2006--2013}{Geomatics Engineer}{\href{http://en.bhidi.com/}{Beijing Engineering Corporation}}{Beijing, China}{Hydropower Engineering}{%
  Delivered geospatial surveying and technical design for large-scale hydropower infrastructure projects in China and overseas.}

\section{Technical Skills}
\cvitem{Domain}{Autonomous-driving operations, SOTA/FOTA lifecycle and release orchestration, fleet platforms, software/map/configuration/calibration delivery, observability and recovery}
\cvitem{Backend \& Data}{Python (Flask, Django), REST APIs, Celery, RabbitMQ, Kafka/Faust, MQTT/EMQX, MySQL, MongoDB}
\cvitem{Frontend}{React, Next.js, Ant Design, HTML/CSS; end-to-end product and workflow design}
\cvitem{Systems}{Linux, Shell, Docker, Git, CI/CD, systemd, rsync, OverlayFS, FRP; deployment, monitoring, failure handling and operational support}
\cvitem{Additional}{Go (Gin), practical C++ and OpenCV, TDD, contract testing, agile delivery}

\section{Education}
\vspace{0.08cm}
\cventry{2002--2006}{Bachelor's Degree in Geomatics Engineering}{\href{https://en.hhu.edu.cn/}{Hohai University}}{Nanjing, China}{}{}

\vspace{0.16cm}

\section{Languages \& Additional Information}
\cvitem{Languages}{English (IELTS Academic 6.5); Japanese (JLPT N2); Mandarin Chinese (native)}
\cvitem{Mobility}{Open to relocation across Germany and to working in international teams}
\cvitem{Interests}{Distance running; certified third-level amateur marathon athlete in Beijing}

\end{document}
